\documentclass[12pt,a4paper]{article}

% ---- Encoding & Font ----
\usepackage{iftex}
\ifPDFTeX
  \usepackage[utf8]{inputenc}
  \usepackage[T1]{fontenc}
\else
  \usepackage{fontspec}  % XeTeX/LuaTeX (tectonic): Unicode nativo (·, —, ¿, ª, §)
\fi
\usepackage[spanish]{babel}
\usepackage{csquotes}

% ---- Page layout ----
\usepackage[top=2.5cm, bottom=2.5cm, left=2.5cm, right=2.5cm]{geometry}
\usepackage{setspace}
\onehalfspacing

% ---- Math ----
\usepackage{amsmath, amssymb, amsthm}
\usepackage{mathtools}
\usepackage{physics}

% ---- Graphics ----
\usepackage{graphicx}
\usepackage{tikz}
\usepackage{float}
\usepackage{caption}

% ---- Tables ----
\usepackage{array}
\usepackage{booktabs}
\usepackage{tabularx}
\usepackage{colortbl}

% ---- Colours ----
\usepackage{xcolor}
\definecolor{notebg}{HTML}{FFF3CD}
\definecolor{noteborder}{HTML}{FFC107}
\definecolor{boxred}{HTML}{F8D7DA}
\definecolor{boxredborder}{HTML}{F5C6CB}

% ---- Boxes ----
\usepackage{mdframed}
\usepackage{tcolorbox}
\tcbuselibrary{most}

\newtcolorbox{keybox}[1][]{
  colback=blue!5,
  colframe=blue!70!black,
  arc=4pt,
  boxrule=1pt,
  fonttitle=\bfseries,
  title={#1}
}

\newtcolorbox{warnbox}[1][]{
  colback=boxred,
  colframe=boxredborder,
  coltitle=black!75,
  arc=4pt,
  boxrule=1pt,
  fonttitle=\bfseries,
  title={#1}
}

\newtcolorbox{notebox}[1][]{
  colback=notebg,
  colframe=noteborder,
  coltitle=black!75,
  arc=4pt,
  boxrule=1pt,
  fonttitle=\bfseries,
  title={#1}
}

% ---- Hyperlinks & TOC ----
\usepackage[hidelinks]{hyperref}
\usepackage{bookmark}

% ---- Enumerate ----
\usepackage{enumitem}

% ---- Pseudocódigo ----
\usepackage[spanish,onelanguage,ruled,vlined]{algorithm2e}
\SetKwInput{KwEntrada}{Entrada}
\SetKwInput{KwSalida}{Salida}

% ---- Miscellaneous ----
\usepackage{icomma}
\usepackage{siunitx}
\usepackage{textcomp}
\usepackage[bottom]{footmisc}

% ---- Title ----
\title{\textbf{Clase 5: Sistemas Lineales, Introducción} \\[0.3em]
        \large{Métodos directos, costo computacional y el problema del pivoteo}}
\author{Transcripto y expandido --- Curso Métodos Numéricos (OpenFING)}
\date{}

\begin{document}

\maketitle
\thispagestyle{empty}
\newpage

\tableofcontents
\newpage

\section{El problema y el enfoque del curso}

Cierra la unidad de errores, punto flotante y redondeo, y arranca la unidad
de \textbf{sistemas de ecuaciones lineales}, el primer problema concreto del
curso. El objetivo es dar un \textbf{enfoque algorítmico}: entender cómo
resolver un sistema grande y cuánto cuesta hacerlo.

El problema formal:
$$Ax = b, \qquad A \in \mathbb{R}^{n\times n}, \quad b \in \mathbb{R}^n$$

$A$ y $b$ son los datos; $x \in \mathbb{R}^n$ es la incógnita.

\begin{keybox}[Existencia y unicidad]
El sistema tiene solución única si y solamente si $\det(A) \neq 0$. El curso
asume desde acá que se trabaja siempre en esa situación (sistema compatible
determinado). Los sistemas rectangulares se posponen para después de los
parciales.
\end{keybox}

\section{Métodos directos vs. métodos iterativos}

\begin{itemize}
  \item \textbf{Método directo}: después de una cantidad finita de pasos,
  produce la solución \textbf{exacta} del problema. Ejemplos: eliminación
  gaussiana, regla de Cramer.
  \item \textbf{Método iterativo}: genera una sucesión de vectores
  $x^{(k)} \in \mathbb{R}^n$ tal que
  $$x^{(k)} \xrightarrow[k \to \infty]{} x^\star$$
  donde $x^\star$ es la solución exacta. Se fija una tolerancia y se detiene
  el cómputo cuando la aproximación se considera suficientemente buena.
\end{itemize}

\begin{notebox}[Motivación de los iterativos]
Un sistema $3\times 3$ es trivial para cualquier método. Un ingeniero se
enfrenta a sistemas mucho más grandes, donde un método directo puede ser
demasiado costoso y conviene algo más rápido que dé una solución aceptable a
menos de una tolerancia. Esta clase y la próxima se dedican íntegramente a
un método \textbf{directo}: la eliminación gaussiana.
\end{notebox}

\section{La regla de Cramer y por qué se descarta}

$$x_i = \frac{\det(A_i)}{\det(A)}$$

donde $A_i$ es la matriz que se obtiene sustituyendo la columna $i$-ésima de
$A$ por el vector $b$.

\subsection{Ejemplo trabajado en clase}

$$\begin{pmatrix} 3 & 1 \\ 1 & -1 \end{pmatrix} \begin{pmatrix} x_1 \\ x_2 \end{pmatrix} = \begin{pmatrix} 5 \\ -1 \end{pmatrix}$$

$$\det(A) = -4, \qquad \det(A_1) = -4, \qquad \det(A_2) = -8$$

$$x_1 = \frac{-4}{-4} = 1, \qquad x_2 = \frac{-8}{-4} = 2$$

\subsection{Flops: la unidad de costo}

Un \textbf{flop} es una operación aritmética de punto flotante (suma, resta,
multiplicación o división). El curso mide costo computacional contando el
\textbf{orden de magnitud} de flops en función de $n$, no el tiempo de reloj
real.

\subsection{Costo de Cramer}

Resolver un sistema $n\times n$ por Cramer requiere $n+1$ determinantes.
Calculando cada uno por desarrollo por filas/columnas, cada uno cuesta del
orden de $n!$ flops (el desarrollo anida determinantes más chicos dentro de
determinantes más grandes):

$$\text{Costo}_{\text{Cramer}} \sim (n+1)\cdot n! = (n+1)!$$

\begin{warnbox}[Cramer escala pésimo]
Implementar Cramer ingenuamente sobre una matriz de $20\times 20$ ``funde la
computadora''. Aun con formas más eficientes de calcular determinantes,
Cramer sigue siendo más caro que la eliminación gaussiana —que ya es cara,
$O(n^3)$ (§\ref{sec:costo-gauss})—. Por eso no se usa en la práctica.
\end{warnbox}

\section{Sistemas triangulares}

\subsection{Definición y determinante}

$A$ es \textbf{triangular superior} si $a_{ij} = 0$ para todo $i > j$. $A$ es
\textbf{triangular inferior} si $a_{ij} = 0$ para todo $j > i$.

\begin{keybox}[Determinante de una matriz triangular]
$$\det(A) = \prod_{i=1}^n a_{ii}$$
En particular, $\det(A) \neq 0 \iff a_{ii} \neq 0$ para todo $i$ — esto
garantiza que la sustitución (§\ref{sec:sustitucion}) nunca divide por cero,
bajo la hipótesis de solución única.
\end{keybox}

\subsection{Sustitución hacia adelante y hacia atrás}
\label{sec:sustitucion}

Los sistemas triangulares se resuelven despejando una incógnita a la vez,
\textbf{secuencialmente}.

\textbf{Sustitución hacia adelante} ($A$ triangular inferior): desde $x_1$
hacia $x_n$.
$$x_1 = \frac{b_1}{a_{11}}, \qquad x_i = \frac{b_i - \displaystyle\sum_{j=1}^{i-1} a_{ij}x_j}{a_{ii}} \quad (i = 2,\dots,n)$$

\textbf{Sustitución hacia atrás} ($A$ triangular superior): el mismo
procedimiento empezando por $x_n$ y avanzando hacia $x_1$.

\subsection{Costo de la sustitución: \texorpdfstring{$O(n^2)$}{O(n\texttwosuperior)}}

Calcular $x_1$ cuesta 1 flop. Calcular $x_i$ ($i>1$) cuesta $2i-1$ flops
($(i-1)$ productos, $(i-1)$ restas, 1 división). Sumando:
$$\sum_{i=1}^{n} (2i-1) = 2\cdot\frac{n(n+1)}{2} - n \sim n^2$$

\begin{notebox}[Truco para estimar $\sum i^\alpha$]
Pensar $\sum_{i=1}^n i^\alpha$ como una suma de Riemann que acota (por arriba
y por abajo) a $\int_0^n x^\alpha\,dx = \frac{n^{\alpha+1}}{\alpha+1}$: la
suma es del orden de $n^{\alpha+1}$. Herramienta general, no específica de
esta cuenta.
\end{notebox}

\begin{keybox}[Costo de la sustitución]
Sustitución hacia adelante o hacia atrás cuesta $O(n^2)$. Duplicar $n$
multiplica el costo por 4.
\end{keybox}

\section{Eliminación gaussiana sin pivoteo}
\label{sec:costo-gauss}

¿Cuánto cuesta \textbf{llevar} un sistema cualquiera a forma triangular?

\subsection{La idea: multiplicadores y combinaciones lineales}

Se trabaja sobre la matriz ampliada $[A \mid b]$, haciendo ceros columna por
columna debajo de la diagonal mediante combinaciones lineales de filas que no
cambian el conjunto solución ni el determinante.

En el paso $k$-ésimo, asumiendo $a_{kk}^{(k)} \neq 0$ (el \textbf{pivote}),
se definen los \textbf{multiplicadores}:
$$l_{ik} = \frac{a_{ik}^{(k)}}{a_{kk}^{(k)}}, \qquad i = k+1,\dots,n$$
y se actualizan las filas $i = k+1,\dots,n$:
$$a_{ij}^{(k+1)} = a_{ij}^{(k)} - l_{ik}\,a_{kj}^{(k)}, \qquad b_i^{(k+1)} = b_i^{(k)} - l_{ik}\,b_k^{(k)}$$
para $j = k+1,\dots,n$. Por construcción $a_{ik}^{(k+1)} = 0$ para todo
$i>k$.

\subsection{Pseudocódigo}

\begin{algorithm}[H]
\KwEntrada{$A \in \mathbb{R}^{n\times n}$, $b \in \mathbb{R}^n$}
\KwSalida{$A$, $b$ transformados: sistema triangular superior equivalente}
\For{$k \gets 1$ \KwTo $n-1$}{
  \For{$i \gets k+1$ \KwTo $n$}{
    $l_{ik} \gets a_{ik} / a_{kk}$\;
    \For{$j \gets k+1$ \KwTo $n$}{
      $a_{ij} \gets a_{ij} - l_{ik}\,a_{kj}$\;
    }
    $b_i \gets b_i - l_{ik}\,b_k$\;
  }
}
\caption{Eliminación gaussiana sin pivoteo}
\end{algorithm}

\subsection{Costo: \texorpdfstring{$O(n^3)$}{O(n\textthreesuperior)}}

Fijado $k$: para cada $i$ (hay $n-k$ valores), calcular $l_{ik}$ cuesta 1
flop, el loop interno en $j$ cuesta $2(n-k)$ flops, más 2 flops para
$b_i$. Total por $i$: $3 + 2(n-k)$, repetido $n-k$ veces:
$$\text{Costo}(k) \sim 3(n-k) + 2(n-k)^2$$
Sumando sobre $k = 1,\dots,n-1$:
$$\sum_{k=1}^{n-1} \left[3(n-k) + 2(n-k)^2\right] \sim \frac{2}{3}n^3$$

\begin{keybox}[Costo total]
$$\text{Costo total (gaussiana + sustitución)} = \frac{2}{3}n^3 + O(n^2) \sim O(n^3)$$
Duplicar $n$ multiplica el costo por 8. El término $n^3$ domina: la
eliminación gaussiana es mucho más cara que la sustitución que le sigue.
\end{keybox}

\section{Pivoteo: qué es y por qué hace falta}

\textbf{Pivotear = intercambiar filas.} En álgebra lineal I, la regla
enseñada era: pivotear \emph{sólo si} el pivote candidato $a_{kk}^{(k)}$ es
exactamente $0$.

\begin{warnbox}[La regla clásica no alcanza en punto flotante]
No basta con pivotear sólo cuando el pivote es exactamente cero. Un pivote
\textbf{cercano a cero} —aunque matemáticamente válido— puede ser igual de
peligroso en una computadora con precisión finita, porque genera
multiplicadores $l_{ik}$ muy grandes que amplifican el error de redondeo.
Se ilustra en §\ref{sec:ejemplo}.
\end{warnbox}

\section{Ejemplo numérico: cuando un pivote chico rompe todo}
\label{sec:ejemplo}

Computadora hipotética de \textbf{5 cifras significativas} (con
truncamiento). La solución exacta del sistema trabajado en clase es
$x = (0,-1,1)$.

Tras el primer paso de eliminación, queda un pivote candidato
$a_{22}^{(2)} = -0{,}01$ — muy cercano a cero pero no exactamente cero.
Siguiendo la regla clásica se lo usa igual, produciendo un multiplicador
enorme:
$$l_{32} = \frac{-2{,}5}{-0{,}01} = -2500$$

Al propagar este multiplicador aparece una resta de números de magnitud muy
distinta ($15\,002$ vs.\ $15\,004{,}5$) que la máquina de 5 cifras trunca en
cada paso, perdiendo precisión. El resultado final es

$$x_3 \approx 0{,}99993, \qquad x_2 \approx -1{,}4, \qquad x_1 \approx -0{,}28$$

muy lejos de $(0,-1,1)$.

\begin{notebox}[El mensaje central]
El error no vino de la computadora ``fallando'', sino de una mala elección
de pivote que la regla clásica de álgebra lineal no detecta. Un pivote
pequeño genera multiplicadores grandes que amplifican el error de redondeo.
\emph{Clase siguiente}: se retoma este mismo ejemplo para mostrar la
estrategia de pivoteo (parcial) que evita el problema — tema de la
descomposición LU.
\end{notebox}

\end{document}
